% The Slide Definitions
%document
\documentclass[10pt]{beamer}
%theme
\usetheme{metropolis}
% packages
\usepackage{color}
\usepackage{listings}
\usepackage[ngerman]{babel}
\usepackage[utf8]{inputenc}
\usepackage{multicol}


% color definitions
\definecolor{mygreen}{rgb}{0,0.6,0}
\definecolor{mygray}{rgb}{0.5,0.5,0.5}
\definecolor{mymauve}{rgb}{0.58,0,0.82}

\lstset{
    backgroundcolor=\color{white},
    % choose the background color;
    % you must add \usepackage{color} or \usepackage{xcolor}
    basicstyle=\footnotesize\ttfamily,
    % the size of the fonts that are used for the code
    breakatwhitespace=false,
    % sets if automatic breaks should only happen at whitespace
    breaklines=true,                 % sets automatic line breaking
    captionpos=b,                    % sets the caption-position to bottom
    commentstyle=\color{mygreen},    % comment style
    % deletekeywords={...},
    % if you want to delete keywords from the given language
    extendedchars=true,
    % lets you use non-ASCII characters;
    % for 8-bits encodings only, does not work with UTF-8
    frame=single,                    % adds a frame around the code
    keepspaces=true,
    % keeps spaces in text,
    % useful for keeping indentation of code
    % (possibly needs columns=flexible)
    keywordstyle=\color{blue},       % keyword style
    % morekeywords={*,...},
    % if you want to add more keywords to the set
    numbers=left,
    % where to put the line-numbers; possible values are (none, left, right)
    numbersep=5pt,
    % how far the line-numbers are from the code
    numberstyle=\tiny\color{mygray},
    % the style that is used for the line-numbers
    rulecolor=\color{black},
    % if not set, the frame-color may be changed on line-breaks
    % within not-black text (e.g. comments (green here))
    stepnumber=1,
    % the step between two line-numbers.
    % If it's 1, each line will be numbered
    stringstyle=\color{mymauve},     % string literal style
    tabsize=4,                       % sets default tabsize to 4 spaces
    % show the filename of files included with \lstinputlisting;
    % also try caption instead of title
    language = Python,
	showspaces = false,
	showtabs = false,
	showstringspaces = false,
	escapechar = ,
}

\def\ContinueLineNumber{\lstset{firstnumber=last}}
\def\StartLineAt#1{\lstset{firstnumber=#1}}
\let\numberLineAt\StartLineAt



\newcommand{\codeline}[1]{
	\alert{\texttt{#1}}
}


% Author and Course information
% This Document contains the information about this course.

% Authors of the slides
\author{Felix Döring, Felix Wittwer, Anton Obersteiner}

% Name of the Course
\institute{Python-Kurs}

% Fancy Logo 
\titlegraphic{\hfill\includegraphics[height=1.25cm]{../templates/fsr_logo_cropped}}



% Custom Bindings
% \newcommand{\codeline}[1]{
%	\alert{\texttt{#1}}
%}


% Presentation title
% TODO Change the topic of the lesson
\title{Regul\"are Ausdr\"ucke}
\date{\today}


\begin{document}

\maketitle

\begin{frame}{Gliederung}
        \begin{multicols}{2}
                \setbeamertemplate{section in toc}[sections numbered]
                \tableofcontents
        \end{multicols}
\end{frame}

\section{Grundlagen}
\begin{frame}{Grundlagen}
	Das \alert{re} Modul der Python Standard Library ist die Python Implementierung von regulären Ausdrücken.\\[.25cm]
	Reguläre Ausdrücke werden verwendet um die Struktur von Text/Sprachen zu beschreiben.
\end{frame}

\begin{frame}{Grundlagen}
	Ein erstellter regul\"arer Ausdruck kann verwendet werden um:
	\begin{itemize}
		\item die Struktur eines Textes zu \"uberpr\"ufen
		\item bestimmte Teile eines Textes zu extraktieren
	\end{itemize}
\end{frame}

\begin{frame}{Grundlagen}
	Die Anwendung eines regulären Ausdruckes nennt man Matching.\\[.25cm]
Matcht der reguläre Ausdruck einem String, bedeutet dies, dass der String die Struktur hat, die der reguläre Ausdruck beschreibt.\\[.25cm]
Reguläre Ausdrücke werden oft auch 'regex' genannt (kurz für \alert{regular expression}).
\end{frame}

\section{Matching Regeln}
\begin{frame}{Matching Regeln}
	Jeder Buchstabe und jede Zahl matcht immer \alert{einmal} sich selbst.\\[.25cm]
	Das bedeutet:\\
	\begin{itemize}
		\item 	Der String \texttt{'a'} matcht der regex \texttt{'a'}.\\
		\item Der String \texttt{'abc'} matcht der regex \texttt{'abc'}, nicht aber regex \texttt{'a'}.\\
		\item \texttt{'4'} matcht \texttt{'4'}, nicht aber \texttt{'5'}, \texttt{'a'} oder \texttt{'41'} usw.
	\end{itemize}
\end{frame}

\subsection{Sonderzeichen}
\begin{frame}{Sonderzeichen}
	\texttt{.} (Punkt) matcht \alert{einem} beliebingen Schriftzeichen, außer \texttt{'\textbackslash{}n'} (newline)\\[.25cm]
	Für die regex \texttt{'.'} gilt:\\
	\begin{itemize}
		\item \texttt{'a'} matcht
		\item \texttt{'b'} matcht
		\item \texttt{'4'} matcht
		\item \texttt{'ab'} matcht nicht, denn \texttt{'.'} ist nur ein Zeichen und \texttt{'ab'} sind zwei.
	\end{itemize}
\end{frame}

\begin{frame}{Sonderzeichen}
	\texttt{[]} matcht jedem der in den Klammern stehenden Zeichen, jedoch nur \alert{einmal} (wie bei \texttt{.}).\\[.25cm]
	Für die regex \texttt{'[abg]'} gilt also:\\
	\begin{itemize}
		\item \texttt{'a'} matcht
		\item \texttt{'b'} matcht
		\item \texttt{'g'} matcht
		\item \texttt{'ab'} matcht nicht (zwei zeichen matchen nicht einem).
	\end{itemize}
\end{frame}

\begin{frame}{Sonderzeichen}
	\texttt{()} erstellt eine Gruppe. Alles was in den Klammern steht, muss genau so vor kommen.\\[.25cm]
	Für die regex \texttt{'(abc)'} gilt:\\
	\begin{itemize}
		\item \texttt{'a'} matcht nicht
		\item \texttt{'ab'} matcht nicht
		\item \texttt{'abc'} matcht
	\end{itemize}
\end{frame}

\begin{frame}{Sonderzeichen}
	\texttt{\textbackslash} escaped ein Sonderzeichen.\\[.25cm]
	Alle hier aufgeführten Sonderzeichen können nicht in einem Pattern vorkommen, als das, was sie eigentlich bedeuten, dafür müssen sie extra markiert werden.\\
	\begin{itemize}
		\item \texttt{'\textbackslash\textbackslash'} als Pattern matcht auf den String \texttt{'\textbackslash'}
		\item \texttt{'\textbackslash.'} als Pattern matcht auf den String \texttt{'.'}
	\end{itemize}
\end{frame}

\begin{frame}{Sonderzeichen}
	\texttt{\^} matcht ab dem Anfang eines Strings oder ab jedem \texttt{\textbackslash{}n}\\[.25cm]
	Für die regex \texttt{'\^{}a'} ergibt sich also:\\
	\begin{itemize}
		\item \texttt{'a'} matcht
		\item \texttt{'ba'} matcht nicht, da der falsche Character am Anfang steht.
		\item \texttt{'aba'} matcht
	\end{itemize}	
\end{frame}

\begin{frame}{Sonderzeichen}
	\texttt{\$} matcht dem Ende eines Strings (oder dem Zeilenende)\\[.25cm]
	Für die regex \texttt{'a\$'} folgt daraus:\\
	\begin{itemize}
		\item \texttt{'a'} matcht
		\item \texttt{'ba'} matcht
		\item \texttt{'bab'} matcht nicht, da der falsche Character am Ende steht
		\item \texttt{'aba'} matcht
	\end{itemize}
\end{frame}

\begin{frame}{Sonderzeichen}
	\texttt{|} ist ein \textbf{ODER}. Entweder die regex davor oder danach muss matchen.\\[.25cm]
	Für die regex `'a|b'` gilt:\\
	\begin{itemize}
		\item \texttt{'a'} matcht
		\item \texttt{'b'} matcht
		\item \texttt{'ab'} matcht nicht, da \texttt{'a|b'} mit [ab] gleichzusetzen ist
	\end{itemize}
\end{frame}

\subsection{Zusammengesetzte Regex}
\begin{frame}{Zusammengesetzte Regex}
	Regular Expressions setzen sich aus kleineren Regular Expressions zusammen.\\
	So kann man z.B. auch festlegen, wie häufig ein Zeichen auftreten soll.
\end{frame}

\begin{frame}{Zusammengesetzte Regex}
	\texttt{*} - Die \alert{vorangestellte} Regex muss 0 - n Mal vorkommen.\\[.25cm]
	Für die regex \texttt{'a*'} gilt:\\
	\begin{itemize}
		\item \texttt{''} matcht
		\item \texttt{'a'} matcht
		\item \texttt{'aa'} matcht
		\item \texttt{'aaaab'} matcht nicht, da ein anderes Zeichen als 'a' vorkommt
	\end{itemize}
\end{frame}

\begin{frame}{Zusammengesetzte Regex}
	\texttt{+} - Die \alert{vorangestellte} Regex muss 1 - n Mal vorkommen.\\[.25cm]
	Für die regex \texttt{'a+'} gilt:\\
	\begin{itemize}
		\item \texttt{''} matcht nicht, da es kein mal vorkommt
		\item \texttt{'a'} matcht
		\item \texttt{'aa'} matcht
		\item \texttt{'ab'} matcht nicht, da ein anderes Zeichen als 'a' vorkommt
	\end{itemize}
\end{frame}

\begin{frame}{Zusammengesetzte Regex}
	\texttt{?} - Die \alert{vorangestellte} regex muss 0 - 1 Mal vorkommen.\\[.25cm]
	Für die regex \texttt{'a?'} gilt:\\
	\begin{itemize}
		\item \texttt{''} matcht
		\item \texttt{'a'} matcht
		\item \texttt{'aa'} matcht nicht, da das Zeichen öfter, als nur ein mal vorkommt
	\end{itemize}
\end{frame}

\begin{frame}{Zusammengesetzte Regex}
	\texttt{\{m\}} - Die \alert{vorangestellte} regex muss genau m Mal vorkommen.\\[.25cm]
	Für die regex \texttt{'y\{3\}'} gilt:\\
	\begin{itemize}
		\item \texttt{'yyy'} matcht
		\item \texttt{'y'} matcht nicht, da es zu wenige Zeichen sind
		\item \texttt{'yyyy'} matcht nicht, da es mehr Zeichen sind
	\end{itemize}
\end{frame}

\begin{frame}{Zusammengesetzte Regex}
	\texttt{\{m,n\}} -  Die \alert{vorangestellte} regex muss m - n Mal vorkommen.\\[.25cm]
	Für die regex \texttt{'y\{2,5\}'} gilt:\\
	\begin{itemize}
		\item \texttt{'yyy'} matcht
		\item \texttt{'y'} matcht nicht, da es zu wenige Zeichen sind
		\item \texttt{'yyyy'} matcht
		\item \texttt{'yyyyyy'} matcht nicht, da es zu viele Zeichen sind
	\end{itemize}
\end{frame}

\subsection{Spezielle Sequenzen}
\begin{frame}{Spezielle Sequenzen}
	Desweiteren gibt es noch spezielle Sequenzen, wie z. B.\\[.25cm]
	\begin{itemize}
		\item \texttt{\textbackslash{}d} für Unicode Ziffern, äquivalent für \texttt{[0-9]}
		\item \texttt{\textbackslash{}D} ist das Gegenteil, alles was keine Unicode Ziffern sind
		\item \texttt{\textbackslash{}s} für alle Whitespace Zeichen, das entspricht \texttt{'[\ \textbackslash{}t\textbackslash{}n\textbackslash{}r\textbackslash{}f\textbackslash{}v]'}
		\item \texttt{\textbackslash{}S} entspricht wieder dem Gegenteil
		\item \texttt{\textbackslash{}w} für alle Unicode Zeichen \texttt{'[a-zA-Z0-9\_]'}
		\item \texttt{\textbackslash{}W} für alle Nicht-Unicode Zeichen
		\item \texttt{[\^{}...]} entspricht allem, was nicht in den Klammern steht
	\end{itemize}
\end{frame}

\section{Methoden}
\begin{frame}[fragile]{Methoden}
	\lstinputlisting[lastline=1]{resources/10_regex/methods.py}
	Wandelt einen String in ein regular expression Objekt um.\\[.25cm]
	\lstinputlisting[firstline=3, lastline=3]{resources/10_regex/methods.py}
	Sucht in `string` nach dem Pattern `pattern`.\\[.25cm]
	\lstinputlisting[firstline=5, lastline=5]{resources/10_regex/methods.py}
	Sucht am Begin des Strings nach dem Pattern.
\end{frame}

\begin{frame}[fragile]{Methoden}
	\lstinputlisting[firstline=7, lastline=7]{resources/10_regex/methods.py}
	Der komplette String und das Pattern müssen übereinstimmen.\\[.25cm]
	\lstinputlisting[firstline=9, lastline=9]{resources/10_regex/methods.py}
	Gibt eine Liste von Strings mit allen passenden Übereinstimmungen zurück.\\[.25cm]
	\lstinputlisting[firstline=11, lastline=11]{resources/10_regex/methods.py}
	Gibt einen Iterator, welcher `match` Objekte beinhaltet zurück\\[.25cm]
	\textbf{\href{https://docs.python.org/3/library/re.html}{Die restlichen Funktionen können in den Docs gefunden werden.}}
\end{frame}

\section{reqular expression Objekt}
\begin{frame}{reqular expression Objekt}
	Ein solches Objekt hat im Großen und Ganzen die selben Methoden, jedoch ohne zusätzliches Pattern, da das Objekt an sich bereits ein Pattern enthält.
\end{frame}

\section{Match Objekt}
\begin{frame}{Match Objekt}
	\lstinputlisting[lastline=1]{resources/10_regex/match.py}
	Gibt die Startposition des Patterns im String zurück.\\[.25cm]
	\lstinputlisting[firstline=3, lastline=3]{resources/10_regex/match.py}
	Gibt die Endposition des Patterns im String zurück.\\[.25cm]
	\lstinputlisting[firstline=5, lastline=5]{resources/10_regex/match.py}
	Gibt ein Tuple zurück \texttt{(m.start([group]), m.end([group]))}
\end{frame}


% nothing to do from here on
\end{document}
